%%%%%%%%%%%%%%%%%%%%%%%%%%%%%%%%%%%%%%%%%
% Professional Newsletter Template
% LaTeX Template
% Version 1.0 (09/03/14)
%
% Created by:
% Bob Kerstetter (https://www.tug.org/texshowcase/) and extensively modified by:
% Vel (vel@latextemplates.com)
%
% 2021.05.03
% zh_TW(LuaLaTeX) version by Edward G.J. Lee <edt1023#gmail.com>
%
% This template has been downloaded from:
% http://www.LaTeXTemplates.com
%
% License:
% CC BY-NC-SA 3.0 (http://creativecommons.org/licenses/by-nc-sa/3.0/)
%
%%%%%%%%%%%%%%%%%%%%%%%%%%%%%%%%%%%%%%%%%

\documentclass[10pt]{article} % The default font size is 10pt; 11pt and 12pt are alternatives

%%%%%%%%%%%%%%%%%%%%%%%%%%%%%%%%%%%%%%%%%
% eBook
% Structural Definitions File
% Version 1.0 (29/12/14)
%
% Created by:
% Vel (vel@latextemplates.com)
%
% This file has been downloaded from:
% http://www.LaTeXTemplates.com
%
% 2021.05.01
% zh_TW(LuaLaTeX) version by Edward G.J. Lee <edt1023@gmail.com>
%
% License:
% CC BY-NC-SA 3.0 (http://creativecommons.org/licenses/by-nc-sa/3.0/)
%
%%%%%%%%%%%%%%%%%%%%%%%%%%%%%%%%%%%%%%%%%

%----------------------------------------------------------------------------------------
%	REQUIRED PACKAGES
%----------------------------------------------------------------------------------------

% for luatex-ja and zh_TW fonts setting.
\usepackage[match]{luatexja-fontspec}
\setmainjfont[BoldFont=Noto Sans CJK TC Medium,
              YokoFeatures = {JFM = zh_TW/quanjiao},
              BoldItalicFont=I.Ngaan,
              ItalicFont=TW-MOE-Std-Kai]{I.MingCP}
\setsansjfont[BoldFont=Noto Sans CJK TC Medium,
              YokoFeatures = {JFM = zh_TW/quanjiao},
              BoldItalicFont=I.Ngaan,
              ItalicFont=TW-MOE-Std-Kai]{Taipei Sans TC Beta Light}
\setmonojfont[BoldFont=Noto Sans CJK TC Medium,
              YokoFeatures = {JFM = zh_TW/quanjiao},
              BoldItalicFont=I.Ngaan,
              ItalicFont=TW-MOE-Std-Kai]{cwTeXFangSong}
\setjfontfamily\iyan{I.Ngaan}

%\usepackage[myfont]{ltj-zhfonts}

\newlength{\zhind}
\settowidth{\zhind}{圖}
\setlength{\parindent}{\zhind}
\renewcommand{\baselinestretch}{1.32}

% 配合 tikz 繪圖，得先由 lua 做繪固指令的 output。 → 似乎新版本已開始直接支援。
%\usepackage{luacode}

% lualatex 無需以下二行，本身即為 UTF-8 based.
%\usepackage[utf8]{inputenc} % Required for inputting international characters
%\usepackage[T1]{fontenc} % Output font encoding for international characters

\usepackage[osf]{libertine} % Use the Libertine font
% 中文文件無需 microtype。
%\usepackage{microtype} % Improves character and word spacing

% 中文章節（只適合 artice/report/book，memoir 要另外設定）。
%\renewcommand{\prepartname}{第}
%\renewcommand{\postpartname}{部}
%\renewcommand{\prechaptername}{第}
%\renewcommand{\postchaptername}{章}
%\renewcommand{\presectionname}{}%  第
%\renewcommand{\postsectionname}{}% 節
\renewcommand{\contentsname}{目次}

\usepackage{tikz} % Required for drawing custom shapes
\definecolor[named]{color01}{rgb}{.2,.4,.6} % Color used in the title page
\usepackage{wallpaper} % Required for setting background images (title page)

\usepackage[unicode=true,bookmarks=true,bookmarksnumbered=false,
            bookmarksopen=false,breaklinks=false,pdfborder={0 0 1},
            backref=section,colorlinks=false]{hyperref} % PDF meta-information specification

%----------------------------------------------------------------------------------------
%	PAPER, MARGIN AND HEADER/FOOTER SIZES
%----------------------------------------------------------------------------------------

\setstocksize{12.5cm}{9.4cm} % Paper size
\settrimmedsize{\stockheight}{\stockwidth}{*} % No trims
\setlrmarginsandblock{18pt}{18pt}{*} % Left/right margins
\setulmarginsandblock{30pt}{36pt}{*} % Top/bottom margins
\setheadfoot{14pt}{12pt} % Header/footer height
\setheaderspaces{*}{8pt}{*} % Extra header space

%----------------------------------------------------------------------------------------
%	FOOTNOTE CUSTOMIZATION
%----------------------------------------------------------------------------------------

%\renewcommand{\foottextfont}{\itshape\footnotesize} % Font settings for footnotes
\setlength{\footmarkwidth}{-.1em} % Space between the footnote number and the text
\setlength{\footmarksep}{.1em} % Space between multiple footnotes on the same page
\renewcommand*{\footnoterule}{} % Remove the rule above the first footnote
\setlength{\skip\footins}{1\onelineskip} % Space between the body text and the footnote

%----------------------------------------------------------------------------------------
%	HEADER AND FOOTER FORMATS
%----------------------------------------------------------------------------------------

\makepagestyle{mio} % Define a new custom page style
\setlength{\headwidth}{\textwidth} % Header the same width as the text
\makeheadrule{mio}{\textwidth}{0.1mm} % Header rule height
\makeoddhead{mio}{\scriptsize{\theauthor\hskip.2cm\vrule\hskip.2cm\itshape{\thetitle}}}{}{} % Header specification
\makeoddfoot{mio}{}{\scriptsize {\thepage \quad \vrule \quad \thelastpage}}{} % Footer specification
\makeoddfoot{plain}{}{\footnotesize {\thepage \quad \vrule \quad \thelastpage}}{} % Pages of chapters
\pagestyle{mio} % Set the page style to the custom style defined above

%----------------------------------------------------------------------------------------
%	PART FORMAT
%----------------------------------------------------------------------------------------

\renewcommand{\partname}{部}
\renewcommand{\partnamefont}{\centering\sffamily\itshape\Huge} % Part name font specification
\renewcommand{\partnumfont}{\sffamily\Huge} % Part number font specification
\renewcommand{\parttitlefont}{\centering\sffamily\scshape} % Part title font specification
\renewcommand{\beforepartskip}{\null\vskip.618\textheight} % Whitespace above the part heading

%----------------------------------------------------------------------------------------
%	CHAPTER FORMAT
%----------------------------------------------------------------------------------------

\makechapterstyle{Tufte}{ % Define a new chapter style
\renewcommand{\chaptername}{第}
\renewcommand{\chapterheadstart}{\null \vskip3.5\onelineskip} % Whitespace before the chapter starts
%\renewcommand{\printchaptername}{\large\itshape\chaptername} % "Chapter" text font specification
\renewcommand{\printchaptername}{\LARGE\itshape\chaptername} % "Chapter" text font specification
\renewcommand{\printchapternum}{\LARGE\kansuji\thechapter\mbox{ }章\\} % Chapter number font specification % dirty method.
\renewcommand{\afterchapternum}{} % Space between the chapter number and text
\renewcommand{\printchaptertitle}[1]{ % Chapter title font specification
\raggedright
\itshape\Huge{##1}}
\renewcommand{\afterchaptertitle}{
\vskip3.5\onelineskip
}}
\chapterstyle{Tufte} % Set the chapter style to the custom style defined above

%----------------------------------------------------------------------------------------
%	SECTION FORMAT
%----------------------------------------------------------------------------------------

\setsecheadstyle{\sethangfrom{\noindent ##1}\raggedright\sffamily\itshape\Large} % Section title font specification
\setbeforesecskip{-.6\onelineskip} % Whitespace before the section
\setaftersecskip{.3\onelineskip} % Whitespace after the section

%----------------------------------------------------------------------------------------
%	SUBSECTION FORMAT
%----------------------------------------------------------------------------------------

\setsubsecheadstyle{\sethangfrom{\noindent  ##1}\raggedright\sffamily\large\itshape} % Subsection title font specification
\setbeforesubsecskip{-.5\onelineskip} % Whitespace before the subsection
\setaftersubsecskip{.2\onelineskip} % Whitespace after the subsection

%----------------------------------------------------------------------------------------
%	SUBSUBSECTION FORMAT
%----------------------------------------------------------------------------------------

\setsubsubsecheadstyle{\sethangfrom{\noindent ##1}\raggedright\sffamily\itshape} % Subsubsection title font specification
\setbeforesubsubsecskip{-.5\onelineskip} % Whitespace before the subsubsection
\setaftersubsubsecskip{.1\onelineskip} % Whitespace after the subsubsection

%----------------------------------------------------------------------------------------
%	CAPTION FORMAT
%----------------------------------------------------------------------------------------

\captiontitlefont{\itshape\footnotesize} % Caption font specification
\captionnamefont{\footnotesize} % "Caption" text font specification

%----------------------------------------------------------------------------------------
%	QUOTATION ENVIRONMENT FORMAT
%----------------------------------------------------------------------------------------

\renewenvironment{quotation}
{\par\leftskip=1em\vskip.5\onelineskip\em}
{\par\vskip.5\onelineskip}

%----------------------------------------------------------------------------------------
%	QUOTE ENVIRONMENT FORMAT
%----------------------------------------------------------------------------------------

\renewenvironment{quote}
{\list{}{\em\leftmargin=1em}\item[]}{\endlist\relax}

%----------------------------------------------------------------------------------------
%	MISCELLANEOUS DOCUMENT SPECIFICATIONS
%----------------------------------------------------------------------------------------

%\setlength{\parindent}{1em} % Paragraph indentation 使用一個中文字縮格

\midsloppy % Fewer overfull lines - used in the memoir class and allows a setting somewhere between \fussy and \sloppy

\checkandfixthelayout % Tell memoir to implement the above
 % Include the document which specifies all packages and structural customizations for this template

\begin{document}

%----------------------------------------------------------------------------------------
%	HEADER IMAGE
%----------------------------------------------------------------------------------------

\begin{figure}[H]
\centering\includegraphics[width=0.3\linewidth]{logo.png}
\end{figure}

%----------------------------------------------------------------------------------------
%	第一頁左邊框
%----------------------------------------------------------------------------------------

\begin{minipage}[t]{.30\linewidth} % Mini page taking up 30% of the actual page
\begin{mdframed}[style=sidebar,frametitle={}] % Sidebar box

%-----------------------------------------------------------

\hypertarget{contents}{\textbf{{\large 重點議題……}}} % \hypertarget provides a label to reference using \hyperlink{label}{link text}
\begin{itemize}
\item \hyperlink{firstnews}{第一則新聞標題} % These link to their appropriate sections in the newsletter
\item \hyperlink{secondnews}{第二則新聞標題}
\item \hyperlink{thirdnews}{第三則新聞標題}
\item \hyperlink{descriptivebox}{短語方框}
\item \hyperlink{quotation}{自定的引言}
\end{itemize}

\centerline {\rule{.75\linewidth}{.25pt}} % Horizontal line

%-----------------------------------------------------------

\textbf{左方框第一標題}

又東三百八十里，曰猨翼之山，其中多怪獸，水多怪魚，多白玉，多腹虫，多怪
蛇，多怪木，不可以上。

\begin{enumerate}
\item 又東三百里，曰堂庭之山，多棪木，多白猿，多水玉，多黃金。
\item 又東三百八十里，曰猨翼之山，其中多怪獸，水多怪魚，多白玉，多腹虫，多怪蛇，多怪木，不可以上。
\item 又東三百七十里，曰瞿父之山，無草木，多金玉。
\item 又東四百里，曰句餘之山，無草木，多金玉。
\end{enumerate}

%-----------------------------------------------------------

\textbf{左方框第二標題}

又東三百七十里，曰杻陽之山，其陽多赤金，其陰多白金。有獸焉，其狀如馬而
白首，其文如虎而赤尾，其音如謠，其名曰鹿蜀，佩之宜子孫。怪水出焉，而東流
注於憲翼之水。其中多玄龜，其狀如龜而鳥首虺尾，其名曰\href{http://www.example.com}{旋龜}，其音如判木，佩
之不聾，可以為底。

\textbf{左方框第三標題}

又東三百里，曰青丘之山，其陽多玉，其陰多青䨼。有獸焉，其狀如狐而九尾，
其音如嬰兒，能食人，食者不蠱。有鳥焉，其狀如鳩，其音若呵，名曰灌灌，佩之不
惑。英水出焉，南流注於即翼之澤。其中多\href{http://www.example.com}{赤鱬}，其狀如魚而人面，其音如鴛鴦，
食之不疥。

%-----------------------------------------------------------
\begin{center}
\captionof*{table}{表格標題名}
\begin{tabular}{llr}
\toprule
\multicolumn{2}{c}{姓名} \\
\cmidrule(r){1-2}
姓 & 名 & 級數 \\
\midrule
周 & 伯通 & $7.5$ \\
李 & 莫愁 & $2$ \\
\bottomrule
\end{tabular}
\end{center}
%-----------------------------------------------------------

\end{mdframed}
\end{minipage}\hfill % End the sidebar mini page
%
%----------------------------------------------------------------------------------------
%	第一頁主版面
%----------------------------------------------------------------------------------------
%
\begin{minipage}[t]{.66\linewidth} % Mini page taking up 66% of the actual page

\hypertarget{firstnews}{\heading{第一則新聞標題}{6pt}} % \hypertarget provides a label to reference using \hyperlink{label}{link text}

又東五百里，曰會稽之山，四方，其上多金玉，其下多砆石。勺水出焉，而南流注於湨。

\begin{center}
\parbox[t]{.70\linewidth}{\textit{又東三百四十里曰堯光之山，其陽多玉，其陰多金。有獸焉，其狀如人而彘鬣，穴居而冬蟄，其名曰猾褢，其音如斲木，見則縣有大繇。}}
\end{center}

\textit{《南次二經》}之首，曰柜山，西臨流黃，北望諸毗，東望長右。英水出焉，西南
流注於赤水，其中多白玉，多丹粟。有獸焉，其狀如豚，有距，其音如狗吠，其名
曰狸力，見則其縣多土功。有鳥焉，其狀如鴟而人手，其音如痺，其名曰鴸，其鳴
自號也，見則其縣多放士。

\begin{wrapfigure}[7]{l}[0pt]{0pt} % In-line figure with text wrapping around it
\includegraphics[width=0.3\textwidth]{placeholder.jpg}
\end{wrapfigure}

東五百里，曰\textit{禱過之山}，其上多金玉，其下多犀、兕，多象。有鳥焉，其狀如鵁，
而白首、三足、人面，其名曰瞿如，其鳴自號也。泿水出焉，而南流注於海。其中
有虎蛟，其狀魚身而蛇尾，其音如鴛鴦，食者不腫，可以已痔。

又東五百里，曰丹穴之山，其上多金玉。丹水出焉，而南流注於渤海。有鳥焉，
其狀如雞，五采而文，名曰鳳皇，首文曰德，翼文曰義，背文曰禮，膺文曰仁，腹
文曰信。是鳥也，飲食自然，自歌自舞，見則天下安寧。

項籍少時，學書不成，去學劍，又不成。項梁怒之。籍曰：「書足以記名姓而已。
劍一人敵，不足學，學萬人敵。」於是項梁乃教籍兵法，籍大喜，略知其意，又不
肯竟學。項梁嘗有櫟陽逮，乃請蘄獄掾曹咎書抵櫟陽獄掾司馬欣，以故事得已。項
梁殺人，與籍避仇於\href{http://www.example.com/}{吳中}。吳中賢士大夫皆出項梁下。每吳中有大繇役及喪，項梁
常為主辦，陰以兵法部勒賓客及子弟，以是知其能。秦始皇帝游會稽，渡浙江，梁
與籍俱觀。籍曰：「彼可取而代也。」梁掩其口，曰：「毋妄言，族矣！」梁以此奇
籍。籍長八尺餘，力能扛鼎，才氣過人，雖吳中子弟皆已憚籍矣。

%-----------------------------------------------------------

\hypertarget{secondnews}{\heading{第二則新聞標題}{6pt}} % \hypertarget provides a label to reference using \hyperlink{label}{link text}

項籍者，下相人也，字羽。初起時，年二十四。其季父項梁，梁父即楚將項燕，
為秦將王翦所戮者也。項氏世世為楚將，封於項，故姓項氏。

\begin{itemize}
\item 又東五百里，曰咸陰之山，無草木，無水。
\item 又東三百五十里，曰羽山，其下多水，其上多雨，無草木，多蝮虫。
\item 又東五百里，曰僕勾之山，其上多金玉，其下多草木，無鳥獸，無水。
\end{itemize}

又東四百里，曰洵山，其陽多金，其陰多玉。有獸焉，其狀如羊而無口，不可殺
也，其名曰䍺。洵水出焉，而南流注於閼之澤，其中多芘蠃。

又東五百里，曰鹿吳之山，上無草木，多金石。澤更之水出焉，而南流注於滂
水。水有獸焉，名曰蠱雕，其狀如雕而有角，其音如嬰兒之音，是食人。

凡《南次二經》之首，自柜山至於漆吳之山，凡十七山，七千二百里。其神狀皆
龍身而鳥首。其祠：毛用一璧瘞，糈用稌。

《南次三經》之首，曰天虞之山，其下多水，不可以上。

又東四百里，曰\href{http://www.example.com/}{令丘之山}，無草木，多火。其南有谷焉，曰中谷，條風自是出。
有鳥焉，其狀如梟，人面四目而有耳，其名曰顒，其鳴自號也，見則天下大旱。

\end{minipage} % End the main body - first page mini page

%----------------------------------------------------------------------------------------
%	第二頁主版面
%----------------------------------------------------------------------------------------

\begin{minipage}[t]{.66\linewidth} % Mini page taking up 66% of the actual page

\hypertarget{thirdnews}{\heading{第三新聞標題}{6pt}} % \hypertarget provides a label to reference using \hyperlink{label}{link text}

\begin{multicols}{2} % Two-column layout

\zhlipsum[6,7,8][name=xiangyu]

\BackToContents % Link back to the contents of the newsletter

\end{multicols}

%----------------------------------------------------------------------------------------
%	短語方框
%----------------------------------------------------------------------------------------

\begin{mdframed}[style=intextbox,frametitle={}] % Sidebar box

\hypertarget{descriptivebox}{\heading{短語方框}{0pt}} % \hypertarget provides a label to reference using \hyperlink{label}{link text}

又東三百七十里，曰侖者之山，其上多金玉，其下多青䨼。有木焉，其狀如穀而
赤理，其汗如漆，其味如飴，食者不飢，可以釋勞，其名曰白䓘，可以血玉。
\begin{enumerate}
\item 又東五百八十里，曰禺槀之山，多怪獸，多大蛇。
\item 右南經之山志，大小凡四十山，萬六千三百八十里。
\item 又東五百里，曰發爽之山，無草木，多水，多白猿。汎水出焉，而南流注於渤海。
\item 又東四百里，至於旄山之尾，其南有谷，曰育遺，多怪鳥，凱風自是出。
\item 又東四百里，至於非山之首，其上多金玉，無水，其下多蝮虫。
\item 又東五百里，曰陽夾之山，無草木，多水。
\item 又東五百里，曰灌湘之山，上多木，無草；多怪鳥，無獸。
\item 又東三百七十里，曰侖者之山，其上多金玉，其下多青䨼。
\end{enumerate}

\BackToContents % Link back to the contents of the newsletter

\end{mdframed}

%----------------------------------------------------------------------------------------
%	QUOTATION
%----------------------------------------------------------------------------------------

\hypertarget{quotation}{\large 自訂引言} % \hypertarget provides a label to reference using \hyperlink{label}{link text}

\begin{quote}
\textsl{『項籍者，下相人也，字羽。初起時，年二十四。其季父項梁，梁父即楚將項燕，
為秦將王翦所戮者也。項氏世世為楚將，封於項，故姓項氏。』} —— \textrm{班固}
\end{quote}

%----------------------------------------------------------------------------------------

\end{minipage}\hfill % End of the main body - second page mini page
\begin{minipage}[t]{.30\linewidth} % Mini page taking up 30% of the actual page

%----------------------------------------------------------------------------------------
%	第二頁右邊框
%----------------------------------------------------------------------------------------

\begin{mdframed}[style=sidebar,frametitle={}] % Sidebar box

\heading{規格}{0pt}

班固：

\begin{enumerate}
\item 楚左尹項伯者，項羽季父也，素善留侯張良。
\item 張良是時從沛公，項伯乃夜馳之沛公軍，私見張良，具告以事，欲呼張良與俱去。
\item 曰：「毋從俱死也。」張良曰：「臣為韓王送沛公，沛公今事有急，亡去不義，不可不語。」良乃入，具告沛公。
\item 沛公大驚，曰：「為之柰何？」張良曰：「誰為大王為此計者？」曰：「鯫生說我曰『距關，
毋內諸侯，秦地可盡王也』。故聽之。

\zhlipsum[10][name=trad]

\end{enumerate}

\BackToContents % Link back to the contents of the newsletter

\end{mdframed}\hfill

%----------------------------------------------------------------------------------------

\centering
\begin{minipage}[t]{.95\linewidth}
\textbf{聯繫資訊：}\\
雙手互搏資訊有限公司\\
黃藥師路，\\
桃花島\\
\href{http://www.example.com}{http://www.ornarejustoultricies.com}\\
\href{http://www.example.com}{http://www.elementumsapien.com}\\
\href{http://www.example.com}{http://www.nuncultrices.com}
\end{minipage}

%----------------------------------------------------------------------------------------

\end{minipage} % End of the sidebar mini page

%----------------------------------------------------------------------------------------

\end{document}
